\documentclass[border=14pt]{standalone}
\usepackage{booktabs}
\usepackage{array}
\usepackage[scaled]{helvet}
\renewcommand\familydefault{\sfdefault}
\newcolumntype{L}[1]{>{\raggedright\arraybackslash}p{#1}}

\begin{document}
\small
\begin{tabular}{@{}L{3.5cm} L{2.7cm} L{3.8cm} L{4.9cm}@{}}
\multicolumn{4}{c}{\normalsize\textbf{Ablation summary: which ``knob'' decides adversarial robustness?}}\\[5pt]
\toprule
\textbf{Factor we varied} & \textbf{Held fixed} & \textbf{Result} & \textbf{What we learned}\\
\midrule
Add vs.\ remove\newline {\footnotesize(FoSR vs.\ ProxyDelete)}
  & over-squashing goal
  & both fail to defend
  & \textbf{direction is not what matters}\\
\addlinespace
Edit criterion\newline {\footnotesize(spectral / similarity / low-rank)}
  & that we edit edges at all
  & only attack-aware editing defends (Jaccard, SVD)
  & \textbf{the criterion is the key}\\
\addlinespace
Edge budget\newline {\footnotesize(253 vs.\ 643 edges)}
  & the method
  & spectral removal still fails
  & it is not a budget artifact\\
\addlinespace
Dataset\newline {\footnotesize(Cora vs.\ Citeseer)}
  & the method
  & same ranking of methods
  & the finding generalizes\\
\addlinespace
Attacker\newline {\footnotesize(non-adaptive vs.\ adaptive)}
  & the defense
  & evades the filter but turns weak
  & potency \& detectability are coupled\\
\addlinespace
Over-squashing measure\newline {\footnotesize(clean vs.\ poisoned graph)}
  & the model
  & decay rate unchanged
  & over-squashing \& fragility are decoupled\\
\bottomrule
\end{tabular}
\end{document}
